\documentclass[aps,prb,onecolumn,nofootinbib,superscriptaddress]{revtex4-2}
\usepackage{amsmath,amssymb,amsthm,mathtools,bm}
\usepackage{booktabs}
\usepackage[colorlinks=true,linkcolor=blue,citecolor=blue,urlcolor=blue]{hyperref}
\usepackage{microtype}

\newcommand{\Dlambda}{\Delta_{\lambda}}
\newcommand{\calC}{\mathcal{C}}
\newcommand{\Nlayer}{N_{\rm layer}}
\newcommand{\pathref}[1]{\path{#1}}

\begin{document}

\title{Worker 12 Note: Lyapunov-Spectrum Evidence in the Class-A Adaptive Fermionic Circuit}
\author{Worker 12}
\date{\today}

\begin{abstract}
I inspected the non-excluded Lyapunov-spectrum evidence, with primary emphasis on
\path{colab_lyapunov/}.  The repository evidence supports a clear qualitative
separation between the uniform topological geometry and the domain-wall geometry.
In the uniform geometry, the finite-time single-particle Lyapunov spectrum is
reported as strictly contracting and separated from zero, with
$\Dlambda=\min_i|\lambda_i|\simeq 0.72$--$0.77$ under perfect correction and
$\simeq 0.95$ under post-selection.  In the domain-wall geometry, the same
diagnostic is suppressed by more than an order of magnitude, decreases over
$N_y=30,40,50$, and is correlated with deviations of the real-space Chern marker
from integer quantization.  The strongest interpretation is that the Lyapunov
calculation detects a slow or marginal single-particle transfer direction tied
to the chiral domain-wall sector.  The principal caveat is that the computation
is finite-time, covers only three circumferences with $T=N_y$, and measures a
single-particle transfer-frame spectrum rather than the full Hermitian
covariance tangent space.
\end{abstract}

\maketitle
\setcounter{tocdepth}{2}
\tableofcontents

\section{Source Scope}

I read \path{REPO_POLICY.md} before inspecting the repository.  I excluded every
path containing \texttt{haining}, \texttt{haoyu}, or \texttt{erroneous}, case
insensitively, and did not read anything under
\path{repo_synthesis/main_pass/worker_notes}, \path{supervisor_notes}, or
\path{boss}.  I did not edit files, execute notebooks, or load large binary
arrays.  The main sources were the Lyapunov writeup
\path{colab_lyapunov/docs/lyapunov_analysis.tex}, the observer/helper script
\path{colab_lyapunov/src/lyapunov_observables_gpu.py}, the Ny-sweep campaign
manifest and run index under \path{colab_lyapunov/gpu_data/lyapunov_spectra/},
and notebook text from the data-generation and characterization notebooks.  I
also noted the older \path{lyapunov_analysis_v2/majorana_choi_transfer_spectrum.ipynb}
as historical transfer-spectrum context, but did not use its cached arrays.

\section{Findings}

\subsection{What Spectrum Is Being Measured}

The Lyapunov spectrum in the Colab campaign is not the full tangent spectrum on
Hermitian covariance perturbations $\delta G\in{\rm Herm}(\Nlayer)$.  The
writeup states that the Benettin frame has columns in the top-layer
single-particle Hilbert space, and that the resulting exponents are those of the
linearized single-particle transfer matrix of the circuit, not of the full
many-body/covariance tangent space.  QR reorthogonalization is performed once
per Markov cycle, and the exponents are normalized by cycle count.  The primary
gap proxy is
\begin{equation}
  \Dlambda(T)=\min_i |\lambda_i(T)| ,
\end{equation}
the final finite-time exponent closest to zero.  This is a stability diagnostic:
an $O(1)$ gap indicates all measured tangent-frame directions contract on a
cycle scale, while a closing $\Dlambda$ indicates a slow, marginal, or
nearly-marginal direction.

\subsection{Uniform Topological Geometry}

For DW off, the evidence is internally consistent.  The writeup reports all
finite-time exponents as strictly negative, with a final-cycle gap
$\Dlambda\simeq0.722(11),0.748(8),0.768(9)$ for $N_y=30,40,50$ under perfect
correction, and $\Dlambda\simeq0.945,0.955,0.952$ under post-selection.  The
document further states that $\Dlambda(t)$ rapidly reaches a plateau within
roughly ten cycles and has no discernible $N_y$ dependence over the available
sizes.  This supports an ergodic-transfer interpretation for the measured
single-particle frame: perturbations are exponentially forgotten and the circuit
flows to the expected Chern-$1$ steady state.

\subsection{Domain-Wall Geometry}

For DW on, the gap is much smaller.  The final-cycle values reported in the
writeup are
\begin{align}
  \Dlambda_{\rm PC}(N_y=30,40,50)&\simeq (0.027,0.024,0.020),\\
  \Dlambda_{\rm PS}(N_y=30,40,50)&\simeq (0.021,0.0076,0.0078).
\end{align}
The perfect-correction series decreases monotonically over the three sizes; the
post-selection values are single-trajectory data and are noisier, but remain
well below the DW-off gaps.  The writeup states that the DW-on gap does not
clearly plateau within $T=N_y$ cycles.  This is important: the evidence points
toward a slow or marginal direction, but the finite-time asymptotic regime is
not definitively established.

\subsection{Chern Marker and Localization Evidence}

The real-space Chern marker is reported to remain quantized at $\calC\simeq 1$
for DW off, while DW on gives sub-integer final values:
$\calC_{\rm PC}\simeq0.978(6),0.954(16),0.954(17)$ and
$\calC_{\rm PS}\simeq0.662,0.863,0.883$ for $N_y=30,40,50$.  The writeup also
reports a sample-by-sample correlation: smaller final Chern marker values tend
to coincide with smaller Lyapunov gaps.  This links the near-zero Lyapunov mode
to the domain-wall sector rather than to a generic QR-frame artifact.  The saved
minimum-$|\lambda|$ vector is described as sharply weighted near the domain-wall
columns $x=5,15$ and extended along $y$, matching a chiral boundary-channel
interpretation.

\section{Interpretation}

The most defensible interpretation is transfer stability rather than a direct
many-body Lyapunov statement.  DW off appears to have a stable, isolated
attracting fixed point in the measured single-particle transfer dynamics.  DW on
has a parametrically slower direction whose gap is suppressed by roughly one to
two orders of magnitude and decreases with circumference.  Because the
domain-wall geometry contains chiral boundary degrees of freedom, the natural
physical assignment is a neutral tangent perturbation that shifts charge or
correlations along the wall.  The Chern-marker correlation and reported
domain-wall localization of the min-gap vector materially strengthen this
assignment.

The evidence does not yet determine the scaling law.  The writeup explicitly
leaves open whether $\Dlambda\sim N_y^{-1}$, another power law, or a logarithmic
closing.  The available data are sufficient for the qualitative conclusion
``DW-on is near-marginal relative to DW-off,'' but insufficient for an exponent
or continuum-limit universality claim.

\section{Evidence Table}

\begin{ruledtabular}
\begin{tabular}{p{0.32\linewidth}p{0.18\linewidth}p{0.42\linewidth}}
Source & Lines & Evidence used \\
\midrule
\pathref{colab_lyapunov/docs/lyapunov_analysis.tex}
& 142--172, 183--204, 206--221
& Defines single-particle tangent propagation, QR accumulation, cycle
normalization, and $\Dlambda=\min_i|\lambda_i|$; states the saved min-gap vector. \\
\pathref{colab_lyapunov/docs/lyapunov_analysis.tex}
& 223--230
& Campaign parameters: $N_x=20$, $N_y=30,40,50$, $T=N_y$, $S=25$ for perfect
correction, $S=1$ post-selection, complex128 A100, canonical GPU entry point. \\
\pathref{colab_lyapunov/docs/lyapunov_analysis.tex}
& 258--269, 292--301
& DW-off spectra are strictly negative and gapped; DW-on spectra develop
near-zero modes and shrink with $N_y$. \\
\pathref{colab_lyapunov/docs/lyapunov_analysis.tex}
& 306--320, 347--355
& Final gap values: DW off $O(0.7$--$0.95)$, DW on $O(0.007$--$0.027)$; scaling
form unresolved. \\
\pathref{colab_lyapunov/docs/lyapunov_analysis.tex}
& 383--403, 464--469
& DW-off Chern marker is quantized; DW-on marker is sub-integer and correlated
sample-by-sample with the Lyapunov gap. \\
\pathref{colab_lyapunov/docs/lyapunov_analysis.tex}
& 478--499, 563--617, 675--688
& Min-gap vector is interpreted as a neutral tangent mode localized at
domain-wall interfaces; protocol dependence and physical interpretation are
discussed. \\
\pathref{colab_lyapunov/src/lyapunov_observables_gpu.py}
& 13--15, 206--215, 267--287, 328--350
& Helper version and canonical entry point; stored arrays include spectra, Chern
marker, charge, min-abs value/index/vector; payload description gives top-layer
complex orbital basis. \\
\pathref{colab_lyapunov/gpu_data/lyapunov_spectra/campaigns/N20_Ny30-50_nsh1_a1-1_S25_cyclesNy/campaign_manifest.json}
& 1--18, 225--228, 229--420
& Manifest records geometry set, channel order, canonical dynamics entry point,
cycles rule, helper version, and actual sample counts. \\
\pathref{colab_lyapunov/gpu_data/lyapunov_spectra/campaigns/N20_Ny30-50_nsh1_a1-1_S25_cyclesNy/run_index.csv}
& 1--13
& Run index confirms all 12 DW/protocol/size cases, $n_{\rm vec}=1200,1600,2000$,
and actual samples: 25 for perfect correction, 1 for post-selection. \\
\pathref{colab_lyapunov/notebooks/data_generation/run_lyapunov_spectra_N20_Ny30-50_S25_cyclesNy.ipynb}
& 33--35, 117--147, 174, 357--404
& Notebook text states the campaign purpose, verifies the
\texttt{lyapunov\_observer} signature, uses canonical
\texttt{run\_markov\_circuit}, and saves spectra/min-vector/Chern/charge outputs. \\
\pathref{lyapunov_analysis_v2/majorana_choi_transfer_spectrum.ipynb}
& 9, 15--20
& Older historical notebook describes Choi/Majorana transfer-spectrum analysis
from cached Lyapunov-related files; useful context but not primary evidence here. \\
\end{tabular}
\end{ruledtabular}

\section{Caveats}

First, the spectra are finite-time spectra with $T=N_y$, and the DW-on gap is
reported not to have cleanly plateaued.  Longer runs such as $T\gg N_y$ are
needed before quoting asymptotic exponents.

Second, the tangent frame lives in the single-particle top-layer orbital space.
This is a meaningful transfer diagnostic, but it is not the complete tangent
spectrum of the nonlinear covariance map on Hermitian matrices.

Third, the size range is narrow.  Three circumferences can show suppression and
a monotone trend, but cannot distinguish $1/N_y$, another power law, or
logarithmic closing.

Fourth, post-selection is deterministic in these runs, with one actual
trajectory.  Its values are useful for protocol comparison but should not be
read as ensemble statistics.

Fifth, the reported localization of the min-gap vector is strong supporting
evidence, but I did not inspect the underlying vector arrays or figures as
binary data.  I relied on the writeup, metadata, and notebook text.

\section{Confidence}

My confidence is high that the repository contains a coherent Lyapunov-spectrum
story: DW off is strongly contracting in the measured transfer frame, while DW
on has a suppressed near-zero direction tied to the domain-wall sector.  My
confidence is moderate, not high, that this already establishes a true
thermodynamic marginal fixed direction.  The finite-time nature, limited sizes,
and single-particle-frame definition are real caveats and should be preserved in
any synthesis.
\end{document}
